\documentclass[10pt,a4paper]{article}
\usepackage[left=1.55cm,right=1.55cm,top=1.25cm,bottom=1.3cm]{geometry}
\usepackage{xeCJK,fontspec}
\usepackage{xcolor}
\usepackage{enumitem}
\usepackage{tabularx}
\usepackage{array}
\usepackage{titlesec}
\usepackage{hyperref}
\usepackage{microtype}
\usepackage{fancyhdr}
\usepackage{lastpage}
\usepackage{fontawesome5}
\usepackage{tikz}

\setmainfont{Times New Roman}
\IfFontExistsTF{SimSun}{\setCJKmainfont{SimSun}}{\setCJKmainfont{FandolSong-Regular}}
\IfFontExistsTF{SimHei}{
  \setCJKsansfont{SimHei}
  \newCJKfontfamily\heiti{SimHei}
}{
  \setCJKsansfont{FandolHei-Regular}
  \newCJKfontfamily\heiti{FandolHei-Regular}
}

\newcommand*{\orcidicon}{%
  \tikz[baseline=-0.55ex]{%
    \fill[fill={rgb,255:red,166;green,206;blue,57}] (0,0) circle (0.72ex);%
    \node[white,font=\sffamily\bfseries,inner sep=0pt,scale=0.38] at (0,0) {iD};%
  }%
}

\definecolor{navy}{HTML}{173F5F}
\definecolor{ink}{HTML}{20242A}
\definecolor{muted}{HTML}{59636E}
\definecolor{rule}{HTML}{9FB3C6}
\definecolor{pale}{HTML}{F2F5F7}
\hypersetup{colorlinks=true,urlcolor=navy,linkcolor=navy}
\color{ink}
\setlength{\parindent}{0pt}
\setlength{\parskip}{0pt}
\linespread{1.06}
\renewcommand{\arraystretch}{1.06}
\setlist[itemize]{leftmargin=1.35em,itemsep=0.10em,topsep=0pt,parsep=0pt,partopsep=0pt,label=\textcolor{navy}{\small$\bullet$}}
\pagenumbering{gobble}

\titleformat{\section}{\heiti\bfseries\color{navy}\fontsize{12.5}{14}\selectfont}{}{0pt}{}[\vspace{0.12em}\color{rule}\titlerule]
\titlespacing*{\section}{0pt}{0.78em}{0.48em}

\newcommand{\daterow}[3]{%
  \noindent
  \begin{minipage}[t]{0.77\textwidth}\vspace{0pt}{\heiti\bfseries #1}\end{minipage}\hfill
  \begin{minipage}[t]{0.21\textwidth}\vspace{0pt}\raggedleft{\color{muted}#2}\end{minipage}\par
  \vspace{0.06em}{\color{muted}#3}\par\vspace{0.16em}
}

\newcommand{\project}[4]{%
  \noindent
  \begin{minipage}[t]{0.79\textwidth}\vspace{0pt}\raggedright{\heiti\bfseries #1}\end{minipage}\hfill
  \begin{minipage}[t]{0.19\textwidth}\vspace{0pt}\raggedleft{\color{muted}#2}\end{minipage}\par
  \vspace{0.18em}{\color{muted}#3}\par
  \vspace{0.08em}{#4}\par\vspace{0.78em}
}

\newcommand{\pub}[2]{\hangindent=1.5em\hangafter=1 {\heiti\bfseries #1}\quad #2\par\vspace{0.28em}}

\pagestyle{fancy}
\fancyhf{}
\renewcommand{\headrulewidth}{0pt}
\renewcommand{\footrulewidth}{0pt}

\begin{document}

% ---------- 紧凑型中文导师初筛页眉 ----------
\begin{tabularx}{\textwidth}{@{}X>{\raggedleft\arraybackslash}p{11.8cm}@{}}
  {\heiti\bfseries\color{navy}\fontsize{24}{26}\selectfont 程荣锋} &
  {\fontsize{8.2}{9.5}\selectfont
  \href{mailto:eb1287945452@mail.scut.edu.cn}{\faIcon{envelope}\ eb1287945452@mail.scut.edu.cn}\quad
  \href{https://crf2004.github.io/zh}{\faIcon{globe}\ crf2004.github.io/zh}\quad
  \href{https://orcid.org/0009-0002-3021-9647}{\orcidicon\ 0009-0002-3021-9647}}\\[0.18em]
  {\color{muted}华南理工大学本科生} &
  {\small\faIcon{map-marker-alt}\ Guangzhou, China}\\[0.08em]
  {\color{muted}预计 2027 年 6 月毕业} &
  {}
\end{tabularx}
\vspace{0.35em}
{\color{navy}\rule{\textwidth}{1.1pt}}

\section{教育背景}
\daterow{华南理工大学｜大数据管理与应用}{2023.09--2027.06}{管理学学士；GPA 3.46/4.0，平均成绩 85.18/100（截至第六学期）}
核心课程：大模型与生成式人工智能（93）、大语言模型与提示工程（93）、机器学习（92）、深度学习（88）。

\section{已获成果与落地实践}
\begin{itemize}
  \item \textbf{论文发表：}参与大模型代码伦理安全评测研究，成果发表于 \textbf{\textit{iScience}}（2026）。
  \item \textbf{知识产权：}第一发明人专利申请 2 项；软件著作权 2 项。
  \item \textbf{临床落地：}联合开发“叮呗心心”心血管健康管理平台，初版已在广东省第二人民医院部署；负责 5 人“叮呗美美”团队的全栈与临床流程开发，产品已完成医院首发。
  \item \textbf{竞赛与项目：}获“互联网+”校赛金奖、广州数学建模高校联赛二等奖、全国大学生数学建模竞赛广东赛区三等奖；连续两年获国家级大学生创新训练计划支持。
\end{itemize}

\section{核心科研经历}
\project{DLEF: Dataset-Label Evidence Fingerprints for Explaining Cross-Dataset Performance in ECG Classification}{2026.06--至今}{第一作者｜指导：张浩彬、程伟彬｜投稿 \textbf{IEEE BIBM 2026}}{
\begin{itemize}
  \item \textbf{问题：}不同中心虽使用相同诊断标签，模型区分阳性与阴性病例所依赖的心电证据仍可能改变；常用性能指标只能显示外部性能下降，难以解释原因。
  \item \textbf{方法：}提出并主导 DLEF 框架：沿节律、传导、QRS 形态、导联特异性 QRS 与复极等证据轴施加受控扰动，以标准化正负贡献差构建数据集--标签指纹，并用方向性距离衡量源端判别证据在目标数据中的保留程度。
  \item \textbf{结果：}在三个公开数据集、六个诊断标签上，DLEF 距离与标签级外部 AUROC 下降密切相关（Spearman $r=0.843$），且优于患病率差异、内部性能与嵌入分布等基线，可在部署前定位高风险标签与迁移方向。
\end{itemize}}

\project{A Matched-Control Perturbation Audit Framework for Evaluating Evidence Use in Deep Learning Models: An Application to 12-Lead ECG Classification}{2026.01--至今}{第一作者｜指导：张浩彬、程伟彬｜投稿 \textbf{IEEE JBHI}}{
\begin{itemize}
  \item \textbf{问题：}模型即使依赖不符合医学知识、或在数据集偏移下不稳定的证据，仍可能获得较高判别性能，因此需要直接审计预测是否由领域相关证据支撑。
  \item \textbf{方法：}提出并主导匹配对照扰动审计：用领域知识定义目标证据及匹配对照，要求目标效应超过最强对照，同时比较扰动算子，并通过冻结权重外部迁移检验证据使用的跨数据集稳定性。
  \item \textbf{结果：}匹配对照推翻了 1AVB 的仅目标区域结论，算子选择改变了 AFIB 与 CRBBB 的判断，外部迁移进一步识别出不稳定的证据依赖；框架由此将证据使用主张划分为可信、受限或不支持。
\end{itemize}}

\section{产品与团队实践}
\project{叮呗心心 / CardioGPT：心血管主动健康大模型}{2025--至今}{联合创始人、核心开发者；2026.08 起任开发负责人}{
负责模型微调、RAG 与全栈实现，将 PPG/ECG、智能问诊和风险评估整合为主动健康管理流程；初版已在广东省第二人民医院部署，现负责 2.0 版本开发。项目连续两年获国家级大学生创新训练计划支持，并于 2026.08 进入公司化运营。}

\project{心肌缺血知识图谱构建与应用}{2025}{华南理工大学学生研究计划｜项目负责人｜指导：董训德}{
带领 5 人跨学科团队，负责问题拆解、团队组织、本体与标注体系设计、质量控制及构建本体维护平台；实现 XGBoost 隐式推理与图谱映射解释分析。}

\newpage
\section{论文与研究成果}
\pub{[1] 已发表}{J. He, S. Qin, H. Zou, et al., \textbf{R. Cheng}, H. Zhang, and W. Cheng. ``CodeSafetyBench: Evaluating and Improving Ethical Safety in Large Language Model Code Generation.'' \textbf{\textit{iScience}}, 2026. \href{https://doi.org/10.1016/j.isci.2026.115073}{doi:10.1016/j.isci.2026.115073}.}
\pub{[2] 已返修}{``Implicit Medical Ethical Deficits in Large Language Models: A Benchmarking and Alignment Study.'' \textbf{\textit{Nature Communications}}；负责审稿回复阶段的主要实验。}
\pub{[3] 第一作者｜在审}{\textbf{R. Cheng et al.} ``DLEF: Dataset-Label Evidence Fingerprints for Explaining Cross-Dataset Performance in ECG Classification.'' \textbf{\textit{IEEE BIBM 2026}}.}
\pub{[4] 第一作者｜在审}{\textbf{R. Cheng et al.} ``A Matched-Control Perturbation Audit Framework for Evaluating Evidence Use in Deep Learning Models: An Application to 12-Lead ECG Classification.'' \textbf{\textit{IEEE JBHI}}.}
\pub{[5] 进行中}{\textbf{PathCLIP: Evidence-Aware ECG--Text Representation Learning.} 面向 ECG--文本对齐与零样本识别，设计匹配数量、层级置乱、残差单独输入与未见标签等对照。}

\section{科研与工程实习}
\project{广东省第二人民医院人工智能研究所｜大模型开发工程师实习生}{2024.11--至今}{医疗数据、知识系统与临床产品研发}{
\begin{itemize}
  \item 基于约 8,000 份心血管病例 PDF 开发标准化患者问答数据流水线，覆盖 LLM 信息抽取、实体消歧、Neo4j 知识表示与混合检索。
  \item 探索将知识图谱路径作为结构支架，生成用于微调、评测及 RAG 的医学问答与推理轨迹数据。
  \item 负责 5 人“叮呗美美”团队的全栈与临床流程开发，产品在医院难治性痤疮诊疗中心成立时完成首发。
\end{itemize}}

\project{广州云展智创信息科技有限公司｜后端与算法开发实习生}{2025.07--2026.01}{医疗数据标准化与健康管理产品}{
面向多机构体检报告设计以 UMLS 为基础的数据模式与标准化流水线，涵盖 LLM 辅助模式演进和医学术语映射；开发 iTwins“健康领航模块”的后端工作流。}

\section{专利、软件著作权与奖项}
\begin{itemize}
  \item \textbf{发明专利：}一种医疗问答输出保障方法、系统及电子设备；第一发明人；申请号 202610582199.6；2026.04 申请，已通过初步审查。
  \item \textbf{发明专利：}一种基于多模态数据融合的痤疮中西医辅助评估方法及系统；第一发明人；2026.07 已提交申请，审查文件待下发。
  \item \textbf{软件著作权：}知识图谱本体智能管理系统 V1.0（2025SR2454864）；基于大模型的心血管主动健康管理智能系统 V1.0（2025SR2080673）。
  \item \textbf{竞赛奖项：}2025 年“互联网+”校赛金奖；2025 年第九届广州数学建模高校联赛二等奖；2025 年全国大学生数学建模竞赛广东赛区三等奖。
\end{itemize}

\section{技术能力}
\begin{tabularx}{\textwidth}{@{}>{\heiti\bfseries}p{3.1cm}X@{}}
编程与框架 & Python、PyTorch、scikit-learn、TypeScript、SQL\\
建模方法 & XGBoost、表示学习、多模态学习、大模型微调、优化\\
大模型与知识系统 & LoRA、DPO、RAG/MQ-RAG、知识图谱、信息抽取、UMLS\\
评估与数据整理 & RAG/检索评估、扰动分析、外部验证、模式设计\\
工程开发 & Neo4j、Next.js、Flask、Git、Docker\\
\end{tabularx}

\end{document}
